% !Mode:: "TeX:UTF-8"
% 任务书中的信息
%% 原始资料及设计要求
\assignReq
{原始资料：国内外相关文献资料 }
{设计技术要求：对于模型未知的学习系统,通过设计离线输入输出测试收集}
{最少的数据,结合迭代学习控制算法,利用采样数据修正系统控制输入,逐步}
{提升对于期望轨迹的跟踪精度,从而完成高精度轨迹跟踪任务,并进行收敛}
{性分析,使用计算机软件实现数值仿真验证。}
%% 工作内容
\assignWork
{针对模型未知的离散线性时不变系统,采用数据驱动的迭代学习控制策略,}
{通过设计离线输入输出测试并采集最少的数据,实现对于给定期望轨迹的良}
{好跟踪,即实际输出与期望轨迹之间的误差单调减小并收敛至零。根据所设}
{计的数据驱动迭代学习控制算法,本文进行了系统跟踪误差的收敛性分析,}
{推导了基于最少输入输出数据的误差收敛条件。文章同时给出了数值仿真,}
{验证了所提数据驱动迭代学习控制方案在轨迹跟踪问题上的有效性。}
%% 参考文献
\assignRef
{[1] Camlibel M K, van Waarde H J, Rapisarda P. The shortest experiment for}
{linear system identification[J]. Systems and Control Letters, 2025, 197: 106045.}
{[2] Bristow D A, Tharayil M, Alleyne A G. A survey of iterative learning}
{control[J]. IEEE Control Systems Magazine, 2006, 26(3): 96-114.}
{[3] Jiang Z. Iterative learning control: A data-driven approach[D]. University}
{of Southampton, 2025.}
{[4] Amann N, Owens D H, Rogers E. Iterative learning control for discrete-time}
{systems with exponential rate of convergence[J]. IEE Proceedings -- Control}
{Theory and Applications, 1996, 143(2): 217-224.}